Bulk fields act on boundary observables through products, brackets and
higher homotopies. Reconstructing the bulk from that action requires
these operations together with their support, cyclic pairing and
quantum state. We study the comparison through explicit polar
reductions and Hamiltonian matrix probes.

On a sphere with a marked double pole, a first-jet Neumann condition
leaves a two-dimensional cyclic cohomology algebra. We construct its
continuous cyclic contraction and the nonlinear map from matrix
connections to residual observables. The full classical reduction
retains gauge stabilizers and boundary transgression terms.

Finite cyclic reductions carry higher operations. Paired smooth cutoffs and
positive-scale heat tensors give explicit quantum contractions and
pole contacts. Normalized finite BV elimination composes through
connected cumulants. Smooth Gaussian responses converge uniformly on
compact sets of matrix backgrounds. Their regularized determinants
and supported normal homotopies have specified coefficient domains.

Hamiltonian probes require their own coefficient complexes.
Semi-free resolutions retain the commutativity relation before passing
to matrices. Coherent derivation actions lift Hamiltonian vector fields.
Central rank, negative cohomology and quantum
torsion distinguish the resulting recipients. Continuous, cyclic and
supported comparisons must preserve these data.

At rank two, a framed
cyclic chart admits an integral quantum Hamiltonian realization.
Global compatibility and formal extension in the framing parameter
have separate, explicit obstructions. The coordinate comparison with
polyvectors precedes the Hochschild comparison: ordinary
antisymmetrization applies on the zero-Poisson fiber, and a nonzero
Poisson tensor requires twisted formality.

Independently prepared compact gauge and Gaussian states have exact
sewing formulas. A scalar interacting preparation exists at fixed
positive smoothing and is unitarily identified with a scalar Fock
interaction on its full form domain. Its cutoff limit converges in
trace norm. The identification of these preparations with the native quantum
pushforward, the interacting continuum master equation and the full
sourced closed-field comparison remain open. The required comparison
must preserve observable insertions and cutting, beyond equality of
partition functions or formal coordinate algebras.
